% !PRINTED 353
\weberpage{353}
%<<CONT
functions of order $\nu$ with respect to $y$, and of order $(4-\nu)$ with respect to $x$. Namely, we have

\weq{4}{
\begin{aligned}
P_0(x,y) &= f(x) \\
P_1(x,y) &= \tfrac{1}{4}\sum_i y_i f'(x_i) \\
P_2(x,y) &= \tfrac{1}{12}\sum_i y_i y_k f''(x_i x_k) \\
P_3(x,y) &= \tfrac{1}{4}\sum x_i f'(y_i) \\
P_4(x,y) &= f(y).
\end{aligned}
}

If $x$ lies on the curve $f$, then $P_0=0$, and one root of equation (2) vanishes. If we further take $(y)$ on the tangent at the point $(x)$, then $P_1(x,y)=0$ as well, and two roots of (2) vanish. The two remaining ones are determined, by (3), through the quadratic equation
\weq{5}{6P_2+4tP_3+t^2P_4=0}
If this equation has two equal roots, the line $(x,y)$ will have yet a second point of contact with the curve $f$, i.\,e. it will be a double tangent. But the condition for this is that the discriminant of equation (5) vanish, or that
\weq{6}{R(x,y) = 2P_3^2 - 3P_4P_2}
be equal to zero. Here $R(x,y)$ is a function homogeneous with respect to each of the two systems $(x)$ and $(y)$, and indeed of degree $2$ for the $x$, of degree $6$ for the $y$. But it is still to be remarked that the equation $R=0$ is also satisfied when $(x)$ is an arbitrary point of the curve, and $(y)$ coincides with $(x)$, because then

\[
P_2(x,x)=P_3(x,x)=P_4(x,x)=f(x)=0
\]

holds. Otherwise, however, $R$ will vanish only when the line $(x,y)$ is a double tangent. We therefore set up the theorem:

1. \gesperrt{The equation $R(x,y)=0$ is, when $(x)$ is a point of the curve $f$ and $(y)$ is a point of the tangent at $(x)$ different from $(x)$, the necessary and sufficient condition for $(x)$ to be a point of contact of a double tangent.}

It is now a matter of deriving from this condition another which contains only the $x$ alone, and which for no
%CONT>>
